\selectlanguage{american}%

\section{Volume Computation by Annealing}

For volume computation, we can apply annealing as follows. We consider
the following estimator that uses a random sample from the distribution
with density proportional to $f_{i}$:

\[
Y=\frac{f_{i+1}(X)}{f_{i}(X)}.
\]
We see that 
\[
\E_{f_{i}}(Y)=\frac{\int f_{i+1}}{\int f_{i}}.
\]
In the DFK algorithm, this ratio is bounded by a constant (in fact
$2$) in each phase, giving a total of roughly $n$ phases since the
ratio of final to initial integrals is exponential. Instead of uniform
densities, we consider 
\[
f_{i}(x)=\exp(-a_{i}\norm x)\chi_{K}(x)\mbox{ or }f_{i}(x)=\exp(-a_{i}\norm x^{2}/2)\chi_{K}(x).
\]
The first form is a useful radial model, while the second is the Gaussian
cooling schedule used below for volume. The variance calculation below
is stated in the more general logconcave-power form, which includes
the Gaussian schedule by taking $f(x)=e^{-\norm x^{2}/2}\chi_{K}(x)$.
The coefficient $a_{i}$ (inverse ``temperature'') will be changed
by a factor of $(1+\frac{1}{\sqrt{n}})$ in each phase, which implies
that $m=\widetilde{O}(\sqrt{n})$ phases suffice to reach the target
distribution. This is perhaps surprising since the ratio of the initial
integral to the final is typically $n^{\Omega(n)}$. Yet the algorithm
uses only $\widetilde{O}(\sqrt{n})$ phases, and hence estimates a
ratio of $n^{\tilde{\Omega}(\sqrt{n})}$ in one or more phases. In
the algorithm below, we assume that $B^{n}\subseteq K\subseteq RB^{n}$
and $f(x)=\exp(-\|x\|^{2}/2)\chi_{K}(x).$

\begin{algorithm2e}[H]

\caption{$\mathtt{LV\ Volume}$}

\SetAlgoLined
\begin{enumerate}
\item For $i=0,\ldots,m=\left\lceil \log_{1+1/\sqrt{n}}(8R^{2}n/\varepsilon)\right\rceil $,
let $a_{i}=2n/(1+\frac{1}{\sqrt{n}})^{i}$ and $f_{i}(x)=f(x)^{a_{i}}.$
\item For $i=0,\ldots,m-1$, 
\begin{enumerate}
\item sample $X^{(1)},X^{(2)},\ldots,X^{(N)}$ from $\pi_{f_{i}}$.
\item Set $Y_{i}=\frac{1}{N}\sum^{N}_{j=1}\frac{f_{i+1}(X^{(j)})}{f_{i}(X^{(j)})}$ 
\end{enumerate}
\item Let $Y=\prod^{m-1}_{i=0}Y_{i}$ and return $Y\cdot(\pi/n)^{n/2}$.
\end{enumerate}
\end{algorithm2e}

The key insight is that even though the expected ratio might be large,
its variance is not.
\begin{lem}
\label{lem:chebychev_ratio} For $X\sim\pi_{f_{i}}$ with $f_{i}(x)=e^{-a_{i}\norm x}\chi_{K}(x)$
for a convex body $K$, or $f_{i}(x)=f(x)^{a_{i}}$ for an integrable
logconcave function $f$, the estimator $Y=\frac{f_{i+1}(X)}{f_{i}(X)}$
satisfies 
\[
\frac{\E\left(Y^{2}\right)}{\E\left(Y\right)^{2}}\le\left(\frac{a^{2}_{i+1}}{(2a_{i+1}-a_{i})a_{i}}\right)^{n}
\]
which is bounded by a constant for $a_{i}=a_{i+1}\left(1+\frac{1}{\sqrt{n}}\right)$.
\end{lem}

\begin{proof}
Let $F(a)=\int_{K}f(x)^{a}\,dx$. Then
\[
\E(Y)=\frac{F(a_{i+1})}{F(a_{i})}\ \mbox{ and }\ \E(Y^{2})=\int\frac{f(x)^{2a_{i+1}}}{f(x)^{2a_{i}}}\cdot\frac{f(x)^{a_{i}}}{F(a_{i})}=\frac{F(2a_{i+1}-a_{i})}{F(a_{i})}.
\]
\begin{align*}
\frac{\E\left(Y^{2}\right)}{\E\left(Y\right)^{2}} & =\frac{F(2a_{i+1}-a_{i})F(a_{i})}{F(a_{i+1})^{2}}.
\end{align*}

By Lemma \ref{lem:ratio_var}, $Z(a)=a^{n}F(a)$ is logconcave in
$a$ and hence
\[
\frac{F(2a_{i+1}-a_{i})F(a_{i})}{F(a_{i+1})^{2}}=\frac{Z(2a_{i+1}-a_{i})Z(a_{i})}{Z(a_{i+1})^{2}}\left(\frac{a^{2}_{i+1}}{(2a_{i+1}-a_{i})a_{i}}\right)^{n}\le\left(\frac{a^{2}_{i+1}}{(2a_{i+1}-a_{i})a_{i}}\right)^{n}.
\]
For $a_{i}=a_{i+1}(1+\frac{1}{\sqrt{n}})$, we have 
\[
\frac{\E\left(Y^{2}\right)}{\E\left(Y\right)^{2}}\le\left(\frac{1/(1+\frac{1}{\sqrt{n}})^{2}}{(\frac{2}{1+\frac{1}{\sqrt{n}}}-1)}\right)^{n}=\left(\frac{1}{(1+\frac{1}{\sqrt{n}})(1-\frac{1}{\sqrt{n}})}\right)^{n}=\left(1+\frac{1}{n-1}\right)^{n}\le4
\]
 for $n\ge2.$ 
\end{proof}
\begin{thm}
Let $r=\E(Y)$ in the LV algorithm. With probability at least $3/4$,
we have $\left|Y-r\right|\le\frac{\varepsilon}{2}r$. 
\end{thm}

\begin{proof}
By Lemmas \ref{lem:prod_var} and \ref{lem:chebychev_ratio}, for
$N\geq100m/\varepsilon^{2}$ we have 
\begin{align*}
\frac{\Var(Y)}{\E(Y)^{2}} & =\prod^{m-1}_{i=0}\left(1+\frac{\Var(Y_{i})}{\E(Y_{i})^{2}}\right)-1\\
 & \le\left(1+\frac{\varepsilon^{2}}{32m}\right)^{m}-1\\
 & \le\frac{\varepsilon^{2}}{16}.
\end{align*}
Now the result follows from Chebyshev's inequality. 
\end{proof}
We can now state the main result of this section.
\begin{thm}[\cite{LV2003}]
The LV algorithm estimates the volume of a convex body in $\R^{n}$
given by a membership oracle to relative error $\varepsilon$ with
probability at least $3/4$ using $\widetilde{O}(n^{4}/\varepsilon^{2})$
oracle queries and $\widetilde{O}(n^{2})$ arithmetic operations per
query.
\end{thm}

\begin{proof}
We first note that since $f_{0}$ corresponds to a Gaussian of variance
$1/2n$, we have 
\[
\int_{B^{n}}f_{0}(x)\ge(1-2^{-n})\int_{\R^{n}}f_{0}(x)\,dx=(1-2^{-n})\left(\frac{\pi}{n}\right)^{n/2}.
\]
Next, by the definition of $m$, $(1+\frac{1}{\sqrt{n}})^{m}\ge8R^{2}n/\varepsilon$,
so $a_{m}=2n/(1+\frac{1}{\sqrt{n}})^{m}\le\varepsilon/(4R^{2})$ and
hence 
\[
f_{m}(x)=e^{-\frac{a_{m}}{2}\norm x^{2}}\ge e^{-\varepsilon/4}.
\]
Together with the previous theorem, this shows that the algorithm's
output is a relative $\varepsilon$ approximation of the volume of
the input convex body with probability at least $3/4$. 

For the complexity, the algorithm has $m$ phases with $O(m/\varepsilon^{2})$
samples per phase. The number of membership queries per sample is
$\widetilde{O}(n^{2}R^{2})$ from an $L_{2}$-warm start using the
hit-and-run algorithm for sampling. By Lemma \ref{thm:L2_warm_start},
samples from phase $i$ provide an $\ell_{2}$-warm start for phase
$(i+1)$. This gives an overall complexity of $O^{*}(n^{3}R^{2})$.
By Theorem \ref{thm:convex_body_rounding}, we can transform the body
to near-isotropic position so that $R^{2}=O(n)$, establishing the
theorem. 
\end{proof}
The LV algorithm has two parts. In the first it finds a transformation
that puts the body in near-isotropic position. The complexity of this
part is $\widetilde{O}(n^{4})$. In the second part, it runs the annealing
schedule, while maintaining that the distribution being sampled is
\emph{well-rounded}, a weaker condition than isotropy. Well-roundedness
requires that a level set of measure $\frac{1}{8}$ contain a constant-radius
ball and that the trace of the covariance (expected squared distance
of a random point from the mean) be bounded by $O(n)$, so that $R/r$
is effectively $O(\sqrt{n})$. To achieve the complexity guarantee
for the second phase, it suffices to use the KLS bound of $\psi_{p}\gtrsim n^{-\frac{1}{2}}$.
Connecting improvements in the Cheeger constant (which is the minimum
isoperimetric ratio over all subsets) directly to the complexity of
volume computation was an open question for a couple of decades. To
apply improvements in the Cheeger constant, one would need to replace
well-roundedness with (near-)isotropy and maintain that. Recent work
shows that a well-rounded convex body can be turned into a near-isotropic
one with complexity $\tilde{O}(n^{3})$. This results in an algorithm
to find a near-isotropic transformation with complexity $O(n^{3.5}\poly\log(nR/r))$. 

A faster volume algorithm is known for well-rounded convex bodies
(any isotropic logconcave density satisfies $\frac{R}{r}=O(\sqrt{n})$
and is well-rounded). This variant of simulated annealing, called
Gaussian cooling, utilizes the fact that the KLS conjecture holds
for a Gaussian density restricted by any convex body, and completely
avoids computing an isotropic transformation.
\begin{thm}[\cite{CV2015}]
The volume of a well-rounded convex body, i.e., with $R/r=O^{*}(\sqrt{n})$,
can be computed using $O^{*}(n^{3})$ oracle calls.
\end{thm}


\section{Integration by Annealing}

The reader will notice that the LV algorithm of the previous section
for volume computation is naturally suited to the more general problem
of integrating logconcave functions. In fact, we can apply it directly,
with minimal changes. The first step of the algorithm is an affine
transformation that we will discuss in the next section. Recall that
$L_{f}(t)$ for a logconcave function $f:\R^{n}\rightarrow\R_{+}$
is the convex level set 
\[
L_{f}(t)=\{x\in\R^{n}:\,f(x)\ge t\}.
\]

\begin{algorithm2e}[H]

\caption{$\mathtt{LV\ Integration}$}

\SetAlgoLined
\begin{enumerate}
\item Put $\pi_{f}$ in near-isotropic position. 
\item Restrict $f$ to the convex body $K=B(0,4\sqrt{n})\cap L_{f}(M_{f}e^{-4n})$.
\item Let $V_{0}$ be an estimate of the volume of $K.$ 
\item For $i=0,\ldots,m$, let $a_{i}=\frac{1}{4n}(1+\frac{1}{\sqrt{n}})^{i},\,a_{m}=1$
and $f_{i}(x)=f(x)^{a_{i}}$, with $m-1$ being the largest integer
such that $a_{m-1}<1$. 
\item For $i=0,\ldots,m-1$, 
\begin{enumerate}
\item sample $X^{(1)},X^{(2)},\ldots,X^{(N)}$ from $\pi_{f_{i}}$.
\item Set $Y_{i}=\frac{1}{N}\sum^{N}_{j=1}f(X^{(j)})^{a_{i+1}-a_{i}}$
\end{enumerate}
\item Let $Y=\prod^{m-1}_{i=0}Y_{i}$ and return $V_{0}\cdot Y$.
\end{enumerate}
\end{algorithm2e}

In words, the algorithm starts with the indicator function over a
convex body containing most of the measure of $f$ and anneals through
a sequence of distributions to reach $f$.
\begin{thm}
The LV algorithm applied to an integrable logconcave function $f:\R^{n}\rightarrow\R_{+}$
estimates the integral of $f$ to relative error $\varepsilon$ with
probability at least $3/4$ using $\widetilde{O}(n^{4}/\varepsilon^{2})$
evaluation oracle queries and $\widetilde{O}(n^{2})$ arithmetic operations
per query.\selectlanguage{english}%
\end{thm}


